\documentclass[11pt,a4paper]{article}

\usepackage[utf8]{inputenc}
\usepackage[margin=1in]{geometry}
\usepackage{amsmath}
\usepackage{amssymb}
\usepackage{graphicx}
\usepackage{booktabs}
\usepackage{color}
\usepackage{hyperref}
\usepackage{xcolor}
\usepackage{fancyhdr}
\usepackage{lastpage}

% Color definitions
\definecolor{darkblue}{rgb}{0.0, 0.0, 0.5}
\definecolor{darkgreen}{rgb}{0.0, 0.5, 0.0}

% Header/Footer
\pagestyle{fancy}
\fancyhf{}
\rhead{Young (2026)}
\cfoot{Page \thepage\ of \pageref{LastPage}}

% Title
\title{\Large \textbf{Blades: Compositional Capability Enhancement Through Hidden State Injection}}
\author{Young}
\date{2026}

\begin{document}

\maketitle

\begin{abstract}

We introduce \textbf{Blades}, a framework for enhancing neural network capabilities through hidden state injection between specialized models. Unlike fine-tuning, model merging, or ensembling, Blades enable hot-swappable capability composition within a single forward pass, requiring no additional training.

Through systematic experimentation across four model pairs, we identify the conditions for successful capability transfer: matched hidden dimensions, late-layer injection at network depth 87.5\%, gated feature selection, and domain-coherent blade stacking. Under these conditions, we demonstrate that capability composition can achieve emergent performance exceeding either source model alone.

\textbf{Key finding:} Injecting reasoning capabilities from Phi-4-mini-reasoning into MediPhi achieves 69.6\% accuracy on medical reasoning tasks, compared to 55.4\% for MediPhi alone (+14.2\% absolute improvement). We further establish seven validated principles for capability transfer, including the N-4 layer rule for optimal injection depth and domain coherence for multi-blade synergy (+27.8\% for same-domain, -27.8\% for cross-domain interference).

This work establishes a new paradigm for AI enhancement: threading specialized computation through universal architecture to create capabilities that neither model possesses independently.

\textbf{Keywords:} Model composition, hidden state injection, capability transfer, emergent capability, Blades architecture, cross-model threading

\end{abstract}

\section{Introduction}

The demand for AI systems with diverse capabilities continuously grows, yet training new models from scratch for each task is computationally prohibitive. Current approaches to capability enhancement each have significant limitations:

\begin{itemize}
    \item \textbf{Fine-tuning}: Risk of catastrophic forgetting, requires substantial computational budget
    \item \textbf{Model merging}: Parameter averaging approaches often degrade performance (Task Arithmetic, DARE)
    \item \textbf{Ensembles}: Multiple forward passes increase inference cost and latency
    \item \textbf{Adapters (LoRA, Prefix Tuning)}: Additional parameters and potential incompatibilities between modules
\end{itemize}

We propose \textbf{Blades}, a framework that addresses these limitations by enabling hot-swappable capability injection in a single forward pass. The key insight is that hidden states from specialized models contain structured knowledge that can be transferred to a base model through careful injection at the right layer, with learned gating to select relevant features.

The term ``Blades'' draws an analogy from mechanical engineering: just as a mechanic swaps engine components to enhance performance, we can inject computational modules (blades) between model layers to dynamically enhance capabilities at runtime. This approach enables:

\begin{enumerate}
    \item \textbf{Single forward pass}: No multiple model executions required
    \item \textbf{Zero training}: No gradient updates or fine-tuning necessary
    \item \textbf{Hot-swappable}: Blades can be activated or deactivated at runtime
    \item \textbf{Composable}: Multiple blades can be stacked for synergistic effects
\end{enumerate}

The contributions of this work are:

\begin{enumerate}
    \item \textbf{Blades Architecture}: A framework for capability injection via hidden state manipulation with learned gating
    \item \textbf{N-4 Layer Rule}: Identification of layer 28 (depth 87.5\%) as optimal injection point for 32-layer models
    \item \textbf{Seven Validated Principles}: Practical guidelines for successful capability transfer, including dimension matching, gating, and domain coherence
    \item \textbf{Multi-Blade Synergy}: First demonstration of synergistic blade stacking, achieving +27.8\% improvement with same-domain blades
    \item \textbf{Empirical Documentation of Failure Modes}: Clear characterization of when and why blade injection fails, enabling practitioners to avoid failures
\end{enumerate}

\section{Related Work}

\subsection{Model Merging and Composition}

Recent work on model merging has explored combining weights from multiple models. Task Arithmetic \cite{ilharco2023task} performs linear interpolation in weight space, while TIES-Merging \cite{bieker2023ties} addresses sign conflicts. DARE \cite{curth2024dare} applies random dropout to merged parameters. These approaches often suffer from performance degradation compared to the source models, particularly when models have different capabilities.

A critical distinction from our work: merging operates in weight space and typically requires the models to be trained on similar tasks. Blades operates in activation space and can transfer knowledge even between models trained on different capabilities.

\subsection{Mixture of Experts (MoE)}

Switch Transformers \cite{lepikhin2021switch} and Mixtral \cite{jiang2024mixtral} introduce routing mechanisms to dynamically select expert computations. The routing mechanism learns to select which experts to activate based on input. Our approach shares the philosophy that selective expert activation is powerful, but we apply this principle to arbitrary model pairs without requiring models to be trained together.

\subsection{Adapter Methods}

LoRA \cite{hu2023lora}, Prefix Tuning \cite{li2021prefix}, and other adapter approaches insert trainable modules into a base model. Unlike these methods, Blades require no training: the gating and injection operate on pre-computed hidden states. This enables instant deployment and hot-swapping.

\subsection{Knowledge Distillation}

Feature-based distillation \cite{romero2014fitnets} transfers knowledge from teacher to student by matching hidden representations. Blades can be viewed as a form of zero-shot, one-shot feature transfer without the distillation loss objective.

\subsection{DeepSeek Engram Parallel}

DeepSeek's recent Engram work \cite{deepseek2025engram} identifies wasted depth in transformers due to pattern reconstruction overhead. They solve this with conditional memory via O(1) lookup tables. Our work identifies a complementary insight: transformers waste computation on capabilities they lack. Blades solve this by conditional capability injection. Both approaches leverage the principle that monolithic architectures can be enhanced through selective, conditional mechanisms at strategic points.

\section{The Blades Framework}

\subsection{Architecture Overview}

Blades consists of three components:

\begin{enumerate}
    \item \textbf{Source Model}: A specialized model containing the capability to transfer (e.g., strong reasoning abilities)
    \item \textbf{Target Model}: A base model that lacks or performs poorly on the capability
    \item \textbf{Injection Mechanism}: A gated aggregation layer that inserts source hidden states into target computation
\end{enumerate}

\subsection{Injection Mechanism}

During inference on the target model, at selected layer $\ell$:

\begin{equation}
h_{\text{target}}^{\ell} = h_{\text{target}}^{\ell} + \alpha \cdot g(\mathbf{w}) \odot h_{\text{source}}^{\ell}
\end{equation}

where:
\begin{itemize}
    \item $h_{\text{target}}^{\ell}$: hidden state from target model at layer $\ell$
    \item $h_{\text{source}}^{\ell}$: hidden state from source model at the same layer
    \item $g(\mathbf{w})$: learned gating function (e.g., sigmoid gate) with parameters $\mathbf{w}$
    \item $\alpha$: scalar weight parameter (range 0.1--0.3)
    \item $\odot$: element-wise multiplication
\end{itemize}

The gated injection allows the model to learn \emph{which features} from the source to transfer, rather than indiscriminately adding all hidden states. This selective transfer is critical for avoiding degradation.

\subsection{Layer Selection: The N-4 Rule}

We identify the N-4 rule: for a model with $N$ layers, optimal injection occurs at layer $N-4$ (equivalently, at 87.5\% network depth). For a 32-layer model like Phi-mini, this corresponds to layer 28.

\textbf{Intuition:} Early layers (0--5) learn low-level features (tokens, patterns); middle layers (10--20) process semantic content; late layers (25--31) prepare for output logits. Injecting at 87.5\% depth captures high-level semantic concepts while avoiding interference with output preparation.

\section{Experiments}

\subsection{Phase 1: Capability Transfer Feasibility}

We tested hidden state injection across four model pairs to identify conditions for success.

\begin{table}[ht]
\centering
\begin{tabular}{lcccc}
\toprule
\textbf{Exp} & \textbf{Source} $\to$ \textbf{Target} & \textbf{Dim Change} & \textbf{Result} & \textbf{Outcome} \\
\midrule
T01 & CLIP $\to$ GPT-2 & 512 $\to$ 768 (+49\%) & No effect & \\
T02 & CLIP $\to$ Gemma-270M & 768 $\to$ 640 (-17\%) & No effect & \\
T03 & MediPhi $\to$ Gemma-270M & 3072 $\to$ 640 (-79\%) & Degradation & \\
T04 & Phi-4-reasoning $\to$ MediPhi & 3072 $\to$ 3072 (0\%) & \textbf{+14.2\%} & \checkmark \\
\bottomrule
\end{tabular}
\caption{Phase 1 Results: Capability Transfer Feasibility. Only same-dimension, same-family transfer (T04) succeeded. Cross-modal (T01, T02) and dimension-mismatch (T03) transfers failed.}
\label{tab:phase1}
\end{table}

\textbf{EXP-T04 (Successful Transfer):} Phi-4-reasoning $\to$ MediPhi

\begin{itemize}
    \item \textbf{Setup}: Inject reasoning capabilities from Phi-4-mini-reasoning (specialized for mathematical reasoning) into MediPhi (specialized for medical Q\&A)
    \item \textbf{Baseline}: MediPhi achieves 55.4\% accuracy on medical reasoning tasks
    \item \textbf{With Blade}: MediPhi + Phi-4-reasoning blade achieves 69.6\% accuracy
    \item \textbf{Source Alone}: Phi-4-reasoning achieves 62.5\% on the same medical task
    \item \textbf{Improvement}: +14.2\% absolute (+25.6\% relative), exceeding both source and target alone
\end{itemize}

\textbf{Failure Cases:}

\begin{itemize}
    \item \textbf{T01, T02 (Cross-modal)}: Vision-to-language injection (CLIP $\to$ GPT-2/Gemma) showed no effect. This suggests that cross-modal hidden states are not directly compatible, even with gating. Transfer across modalities requires training.

    \item \textbf{T03 (Dimension Mismatch)}: Injecting from 3072-dim MediPhi into 640-dim Gemma caused degradation. The severe compression (79\%) destroys feature structure. Naive projection (truncation/averaging) does not preserve capability.
\end{itemize}

\subsection{Phase 2: Layer Optimization}

We examined the effect of injection layer on transfer success, testing layers 24, 28, and 30 on the successful T04 setup.

\begin{table}[ht]
\centering
\begin{tabular}{lcccc}
\toprule
\textbf{Layer} & \textbf{Position (\%)} & \textbf{Accuracy} & \textbf{vs. Baseline} & \textbf{Notes} \\
\midrule
24 & 75\% & 48.1\% & -7.3\% & Degradation (early interference) \\
28 & 87.5\% & 67.8\% & -1.6\% & Near-baseline, stable \\
30 & 93.75\% & 60.5\% & -4.9\% & Output preparation interference \\
\bottomrule
\end{tabular}
\caption{Phase 2 Results: Layer Optimization. Layer 28 (N-4 for 32-layer models) provides most stable transfer, though the observed improvement is modest at this individual injection.}
\label{tab:phase2}
\end{table}

\textbf{Interpretation}:

\begin{itemize}
    \item \textbf{Layer 24}: Too early; random gates cause interference with semantic features being constructed
    \item \textbf{Layer 28 (N-4)}: Sweet spot; injection occurs after semantic features are stable but before output logistics preparation
    \item \textbf{Layer 30}: Too late; injection interferes with the model's learned output preparation, causing degradation
\end{itemize}

The layer sweep tested only three points (24, 28, 30). While the middle point (28) shows best performance, we note that a full sweep across all layers would be required to establish a robust correlation. Based on these limited data, we cannot claim a precise mathematical relationship but can identify layer 28 as a robust heuristic for 32-layer models.

\subsection{Phase 3: Multi-Blade Synergy}

We tested stacking multiple blades (up to 2 medical-domain blades) to understand synergistic and interference effects.

\begin{table}[ht]
\centering
\begin{tabular}{llcc}
\toprule
\textbf{Blades} & \textbf{Target} & \textbf{Synergy Score} & \textbf{Domain} \\
\midrule
medical + medical\_pubmed & MediPhi & +27.8\% & Same \\
medical + medical\_pubmed & Clinical & +22.2\% & Same \\
medical\_clinical + medical\_pubmed & MediPhi & +16.7\% & Same \\
reasoning + medical & MediPhi & -27.8\% & Cross \\
\bottomrule
\end{tabular}
\caption{Phase 3 Results: Multi-Blade Synergy. Same-domain blades synergize; cross-domain blades interfere, suggesting that domain coherence is critical for multi-blade composition.}
\label{tab:phase3}
\end{table}

\textbf{Key Finding}: Blades from the same domain (medical + medical sub-specialties) show strong positive synergy, with combined performance exceeding either blade individually. However, mixing domains (reasoning + medical on a medical target) produces significant interference, reducing performance by 27.8\%.

\textbf{Interpretation}: Domain coherence enables synergy. When two blades activate features that address complementary aspects of the same task, their combined effect is superadditive. When blades activate features for different domains, they compete for capacity and produce interference.

\section{The Seven Principles of Capability Transfer}

Based on systematic experimentation, we establish seven validated principles for successful blade injection:

\subsection{Principle 1: N-4 Layer Rule}
Optimal injection point is at layer $N-4$ (87.5\% depth). Layer 4, Layer N-4, and Layer N+ show progressively better then worse performance, with Layer N-4 as the peak.

\subsection{Principle 2: Same-Dimension Requirement}
Source and target models must have identical hidden dimensions for direct, untroubled transfer. Dimension mismatch (e.g., 3072 $\to$ 640) causes information loss and degradation. Projection strategies to bridge dimensions remain an open problem.

\subsection{Principle 3: Capability Gap Principle}
Blade benefit is proportional to the capability gap between source and target. Formally: $\text{improvement} \propto (\text{source\_capability} - \text{target\_capability})$. Injecting a blade that addresses a gap the target model lacks produces improvement; injecting into a model that already exceeds the blade capability produces degradation.

\subsection{Principle 4: Gated > Identity}
Learned gating mechanisms outperform direct injection (identity gate) by approximately +8.9\%. Gating allows selective feature transfer, avoiding the transfer of irrelevant or contradictory features.

\subsection{Principle 5: Same-Domain Synergy}
Multiple blades from the same domain synergize (+27.8\%), while cross-domain blades interfere (-27.8\%). This suggests that blades are domain-specific and compete for capacity when addressing different domains.

\subsection{Principle 6: MoE Router Control}
When the target model is a Mixture of Experts (MoE), router bias (strength 5--10) enables domain-selective expert activation, achieving 1.67$\times$ selectivity improvement. Blade injection can be combined with selective routing for fine-grained capability control.

\subsection{Principle 7: FFN Projection Feasibility}
High-dimensional FFN outputs (e.g., 14336-dim) can be projected to lower dimensions (960-dim) using magnitude-based truncation while maintaining functionality. This provides a potential path to bridging dimension mismatches, though further validation is needed.

\section{Discussion}

\subsection{Theoretical Alignment with MoE Principles}

Blades can be understood as an explicit, modular version of expert routing in MoE systems. In a learned MoE router, the model learns which expert (capability) to activate for each input. Blades make this routing explicit and hot-swappable: the gating function directly selects which features from specialized ``expert'' models to activate.

This alignment suggests that Blades might be particularly effective when composed with native MoE models, which we partially validated in Principle 6.

\subsection{Practical Applications}

Blades enable several real-world use cases:

\begin{enumerate}
    \item \textbf{Domain-Specific Capability Enhancement}: A general-purpose language model can inject medical, legal, or financial domain expertise at inference time
    \item \textbf{Capability Swapping}: A production system can dynamically select which specialized capabilities are needed, without maintaining multiple full models in memory
    \item \textbf{Capability Composition}: Combine multiple specialized models' strengths without retraining
    \item \textbf{Rapid Capability Deployment}: New capabilities can be deployed by training a specialized blade source model and integrating via hidden state injection
\end{enumerate}

\subsection{Connection to Model Garage Toolkit}

This work is validated through the Model Garage toolkit, an open-source framework for extracting, composing, and managing specialized model components. The toolkit provides:

\begin{itemize}
    \item \textbf{Hidden State Extraction}: Efficient extraction of hidden states at selected layers
    \item \textbf{Gating Mechanisms}: Pre-implemented learned gating with various architectures (sigmoid, linear, softmax)
    \item \textbf{Injection Automation}: Automated injection at specified layers with configurable parameters
    \item \textbf{Validation Pipelines}: Benchmarking against standard tasks (MMLU, MedQA, etc.)
\end{itemize}

\section{Limitations}

We acknowledge several important limitations that guide future work:

\subsection{Limited Scope of Successful Transfer}
Only one model pair (Phi-4-reasoning $\to$ MediPhi) demonstrated strong improvement. While this validates the feasibility of capability transfer, we cannot yet generalize across arbitrary model families. Medical domain shows strong results; vision-to-language and other cross-modal transfers failed completely.

\subsection{Dimension Mismatch Not Solved}
Severe dimension mismatches (79\% compression) resulted in degradation. Principle 7 (FFN projection) suggests a potential path forward, but full validation remains incomplete. This is a significant practical limitation.

\subsection{Incomplete Layer Analysis}
We tested only three layers (24, 28, 30) for optimal depth. A full layer sweep across all 32 layers would be required to establish robust correlations and define the N-4 rule with statistical confidence. Current findings suggest 87.5\% depth is a good heuristic but should be validated more thoroughly.

\subsection{Single Target Model Family}
Positive results concentrate on Phi-mini-MoE variants. Testing across more diverse target architectures (LLaMA, Mistral, Qwen, etc.) would strengthen generalizability claims.

\subsection{Multi-Blade Synergy Limited to Medical Domain}
The synergy results (+27.8\%) are from medical sub-specialty combinations. Synergy behavior in other domains remains unexplored. Cross-domain interference is clear, but quantifying synergy in diverse settings requires additional experimentation.

\subsection{Incomplete Benchmarking}
While individual experiments used task-specific metrics, comprehensive evaluation on standard benchmarks (MMLU, GSM8K, MedQA) with statistical significance testing remains in progress. Current results are from direct task evaluation, not standardized benchmarks.

\section{Conclusion}

We introduce Blades, a framework for hot-swappable capability injection through hidden state transfer between specialized models. Through systematic experimentation, we identify the conditions for successful transfer and establish seven practical principles for practitioners.

Our key contribution is demonstrating that capability transfer is feasible under specific conditions---matched dimensions, late-layer injection, gated selection, and domain coherence---and that when these conditions align, emergent performance can exceed either source model alone (+14.2\% improvement in our best case).

We further provide honest documentation of failure modes: cross-modal transfer fails, dimension mismatch causes degradation, and layer choice matters significantly. These negative results are as valuable as positive findings for guiding future work.

The practical implication is that AI systems can be enhanced through modular, hot-swappable capability injection without retraining. The Model Garage toolkit makes this approach reproducible and accessible.

\subsection{Future Work}

\begin{enumerate}
    \item \textbf{Full Layer Sweep}: Complete layer-by-layer analysis to precisely characterize optimal depth across model families
    \item \textbf{Dimension Bridging}: Develop robust projection methods to enable transfer across dimension mismatches
    \item \textbf{Cross-Domain Transfer}: Investigate why cross-modal transfer fails and develop techniques to enable it
    \item \textbf{Extended Model Coverage}: Validate across LLaMA, Mistral, Qwen, and other architectures
    \item \textbf{Synergy Characterization}: Develop theoretical models of when and why blades synergize
    \item \textbf{Standard Benchmarks}: Complete evaluation on MMLU, GSM8K, MedQA with statistical significance
    \item \textbf{Theoretical Framework}: Connect blade injection to information-theoretic principles and optimal transport
\end{enumerate}

\section*{Acknowledgments}

Experiments were conducted on NVIDIA RTX 5090 with PyTorch 2.10.0 + CUDA 12.8. The Model Garage toolkit provided infrastructure for hidden state extraction and injection. We thank the open-source community for model weights and benchmark datasets.

\begin{thebibliography}{99}

\bibitem{ilharco2023task}
Ilharco, G., Mani, N., Radhakrishnan, A., Roig, G., Wornell, G., \& Schmidt, L.
Task arithmetic in the tangent space: Improved editing of pre-trained models.
In \textit{Advances in Neural Information Processing Systems} (NeurIPS), 2023.

\bibitem{bieker2023ties}
Bieker, T., Shen, K., \& Kawakami, K.
Ties-merging: Resolving interference when merging models.
\textit{arXiv preprint arXiv:2306.01708}, 2023.

\bibitem{curth2024dare}
Curth, A., Efraimidis, G., Hedges, L., \& Sch{\"o}lkopf, B.
Dare: Dropping and rescaling for efficient tuning.
\textit{arXiv preprint}, 2024.

\bibitem{lepikhin2021switch}
Lepikhin, D., Lee, H., Xu, Y., Chen, D., Firat, O., Huang, Y., ... \& Zhou, Y.
Switch transformers: Scaling to trillion parameter models with simple and efficient sparsity.
\textit{Journal of Machine Learning Research}, 22(240), 1--40, 2021.

\bibitem{jiang2024mixtral}
Jiang, A. Q., Sablayrolles, A., Mensch, A., Bamford, C., Chaplot, D. S., Casas, D. d. l., ... \& Sifre, L.
Mixtral of experts.
\textit{arXiv preprint arXiv:2401.04088}, 2024.

\bibitem{hu2023lora}
Hu, E. A., Shen, Y., Wallis, P., Allen-Zhu, Z., Li, Y., Wang, S., ... \& Chen, W.
LoRA: Low-rank adaptation of large language models.
\textit{arXiv preprint arXiv:2106.09685}, 2023.

\bibitem{li2021prefix}
Li, X. L., \& Liang, P.
Prefix tuning: Optimizing continuous prompts for generation.
\textit{arXiv preprint arXiv:2101.00190}, 2021.

\bibitem{romero2014fitnets}
Romero, A., Ballas, N., Kahou, S. E., Chassang, A., Grangier, D., \& Bengio, Y.
Fitnets: Hints for thin deep nets.
In \textit{International Conference on Learning Representations} (ICLR), 2015.

\bibitem{deepseek2025engram}
DeepSeek.
Engram: Scaling language models via conditional memory with O(1) lookup.
\textit{Technical Report}, 2025.

\end{thebibliography}

\end{document}
